\chapter{}
\label{chap:2}

\section{Умова}
Нехай $f \in B_{n}$, $AI (f) = d$ та $g$ -- афінна булева функція вiд $n$ змінних. Довести, що
$d - 1 \leq AI (f \oplus g) \leq d + 1$.

\section{Розв'язання}
\begin{proof}
    Скористаємося нерівністю, яку довели в першій задачі на практиці 4:
    \begin{equation}
        \label{eq:algebraic_immunity_inequality}
        \AI(f \oplus g) \leq \AI(f) + \AI(g) \quad \forall\, f, g \in B_{n}.
    \end{equation}
    Спершу встановимо, чому дорівнює алгебраїчна імунність функції $g$:
    \begin{itemize}
        \item Якщо $g = const$ ($g = 0$ або $g = 1$), то $\AI(g) = 0$: справді, $1 \in \Ann(0)$ і $1 \in \Ann(1 \oplus 1) = \Ann(0)$, 
            тому мінімальний степінь анулятора дорівнює~$0$.
        \item Інакше $\deg g = 1$, і оскільки $g \neq 0$, маємо $g \in \langle g \rangle \setminus \{0\}$, 
            тобто $\min \deg \langle g \rangle \leq 1$. Простими словами: оскільки $g$ -- афінна (а це підвид 
            лінійної, степінь у неї максимум 1). Отже, $\AI(g) \leq 1$.
    \end{itemize}
    \textbf{Верхня оцінка} випливає нерівності~\eqref{eq:algebraic_immunity_inequality}:
    \begin{equation*}
        \AI(f \oplus g) \leq \AI(f) + \AI(g) \leq d + 1.
    \end{equation*}
    \textbf{Нижню оцінку} можна отримати, якщо покласти $f = (f \oplus g) \oplus g$. Застосувавши~\eqref{eq:algebraic_immunity_inequality} 
    до такого запису $f$:
    \begin{equation*}
        \AI(f) \leq \AI(f \oplus g) + \AI(g)
    \end{equation*}
    Підставимо значення:
    \begin{equation*}
        d \leq \AI(f \oplus g) + 1 \Rightarrow \AI(f \oplus g) \geq d - 1
    \end{equation*}
    Об'єднуючи обидві оцінки маємо те, що і треба було показати:
    \begin{equation*}
        d - 1 \leq \AI(f \oplus g) \leq d + 1.
    \end{equation*}
\end{proof}
